\documentclass[11pt,a4paper]{article}
\usepackage[utf8]{inputenc}
\usepackage[T1]{fontenc}
\usepackage{amsmath,amssymb}
\usepackage{graphicx}
\usepackage{hyperref}
\usepackage{booktabs}
\usepackage[margin=1in]{geometry}

\title{Task-Appropriate Hybrid Architectures:\\
A Unified Theory of Gain Across Metrics, Architectures, and Tasks}
\author{Erdem Esa\\ \small Independent Researcher}
\date{\today}

\begin{document}
\maketitle

\begin{abstract}
Hybrid systems do not win unconditionally: their gain is a function of
metric type, architecture type, and task structure. We provide an exact
identity for linear metrics and a per-class extension for macro $F_1$,
which we show is necessary because the linear weight does not collapse
for non-linear metrics. We further show that Kronecker architectures are
$8\times$--$12.6\times$ more parameter-efficient on compositional and
real-image tasks, with the hierarchical variant statistically superior
on image-like inputs ($50$ seeds, $p<0.0001$). These results suggest that hybrid system
design can be made principled: formula-driven for metrics,
structure-driven for tasks.
\end{abstract}


\section{Introduction}
\label{sec:intro}

Hybrid systems that combine symbolic reasoning with neural learning have
become a central theme of modern AI research. Yet the question of
\emph{when} a hybrid system outperforms its neural-only baseline remains
largely heuristic. Prior work has proposed empirical guidelines, but no
unified predictive framework exists.

Our earlier work \cite{esa2026preprint1} established a closed-form
formula for the gain of a hybrid symbolic-neural system over a neural-only
baseline, valid for linear metrics:
\begin{equation}
\text{gain} = \text{cov}_{\text{sym}} \cdot (\text{acc}_s - \text{acc}_n),
\label{eq:linear}
\end{equation}
where $\text{cov}_{\text{sym}}$ is the symbolic arm's coverage and
$\text{acc}_s, \text{acc}_n$ are the symbolic and neural accuracies.
This formula was validated on 12 languages, 4 task types, and 2 metrics
with mean prediction error $0.002$.

The present paper addresses three open questions that the linear formula
does not cover:

\begin{enumerate}
\item \textbf{Metric scope.} Does the formula generalise to non-linear
metrics such as $F_1$?
\item \textbf{Architecture choice.} When is a bilinear Kronecker
architecture preferable to a standard dense MLP?
\item \textbf{Task structure.} How does the gain depend on whether the
task is random-linear or compositional?
\end{enumerate}

\paragraph{Contributions.}
\begin{enumerate}
\item We extend the linear formula to a per-class identity that is
\emph{exact} for macro $F_1$ (Section~\ref{sec:f1}).
\item We show that the Kronecker architecture is $8\times$ more
parameter-efficient on compositional tasks, and that hierarchical
Kronecker statistically dominates dense MLPs on image-like compositional
tasks (Section~\ref{sec:kronecker}).
\item We provide a decision tree that maps task type to architecture
choice, backed by 20--50-seed paired $t$-tests.
\item We document negative results honestly (Sections~\ref{sec:negatives}
and~\ref{sec:discussion}).
\end{enumerate}

\paragraph{Terminology.}
Throughout the paper we distinguish three levels of claim:
\emph{exact identity} (mathematically proven, e.g.,
Equations~\ref{eq:linear-exact}, \ref{eq:f1-macro}),
\emph{predictive formula} (empirically validated approximation,
e.g., Equation~\ref{eq:linear}), and
\emph{hypothesis} (tested but not proven, e.g.,
Section~\ref{sec:kronecker}).


\section{Theory}
\label{sec:theory}

\subsection{Linear Metric Formula}
\label{sec:linear}

Let $s(x) \in \{0, 1, \bot\}$ denote the symbolic arm's output where
$\bot$ is abstention, and $n(x) \in \{0, 1\}$ the neural arm's output.
The hybrid rule is $h(x) = s(x)$ if $s(x) \neq \bot$, else $n(x)$.

Define coverage $\text{cov} = \Pr[s(x) \neq \bot]$ and
$\text{acc}_n^\bot = \Pr[n(x) = y \mid s(x) = \bot]$. For any metric
$M$ that is \emph{linear} in the per-item correctness indicator,
the hybrid's value decomposes exactly as:
\begin{equation}
M_h = \text{cov} \cdot M_s + (1 - \text{cov}) \cdot M_n^\bot.
\label{eq:linear-exact}
\end{equation}
Subtracting $M_n = \text{cov} \cdot M_n + (1 - \text{cov}) \cdot M_n$,
we obtain the gain formula:
\begin{equation}
\text{gain} = \text{cov} \cdot (M_s - M_n) + (1-\text{cov}) \cdot (M_n^\bot - M_n).
\label{eq:linear-gain}
\end{equation}
The simple predictive form drops the second term when
$M_n^\bot \approx M_n$; the residual is bounded by
$\varepsilon = (1-\text{cov}) \cdot |M_n^\bot - M_n|$.


\subsection{F1 Per-Class Formula}
\label{sec:f1}

Macro $F_1$ is not linear in the per-item indicator, so the identity
of Equation~\ref{eq:linear-exact} fails. We derive an \emph{exact
algebraic identity} instead (not a new theorem: it follows directly
from the $F_1$ definition and the hybrid decomposition).

Let $K$ be the number of classes. For class $c$, define the
\emph{positive mass} of a predictor as
\begin{equation}
D^c = 2\,\mathrm{TP}^c + \mathrm{FP}^c + \mathrm{FN}^c.
\label{eq:D-def}
\end{equation}
The classwise $F_1$ is $F_1^c = 2\,\mathrm{TP}^c / D^c$.

For the hybrid system, both TP and D decompose across the answered
and abstained subsets:
\begin{align}
\mathrm{TP}^c_h &= \mathrm{TP}^c_s + \mathrm{TP}^c_n, \\
D^c_h &= D^c_s + D^c_n.
\end{align}

Therefore:
\begin{align}
F_1^c_h &= \frac{2(\mathrm{TP}^c_s + \mathrm{TP}^c_n)}{D^c_s + D^c_n} \\
&= \frac{D^c_s \cdot F_1^c_s + D^c_n \cdot F_1^c_n}{D^c_s + D^c_n} \\
&= w^c \cdot F_1^c_s + (1-w^c) \cdot F_1^c_n,
\label{eq:f1-class}
\end{align}
where the per-class weight is
\begin{equation}
w^c = \frac{D^c_s}{D^c_s + D^c_n}.
\label{eq:w-def}
\end{equation}

Averaging over classes gives the macro identity:
\begin{equation}
F_{1,h} = \frac{1}{K} \sum_{c=1}^{K}
\left[ w^c \, F_1^c_s + (1 - w^c) \, F_1^c_n \right].
\label{eq:f1-macro}
\end{equation}

\paragraph{Relation to linear formula.}
For a \emph{linear} metric, one can show that
$\sum_c D^c = 2N$ holds identically for any prediction, so
the aggregate weight collapses to $w = \text{cov}$. For macro $F_1$,
however, the weights $w^c$ are per-class and do \emph{not} collapse:
class imbalance on the abstained subset breaks the identity.
This is precisely why the linear formula fails for $F_1$ and the
per-class formula succeeds.

\paragraph{Why this matters.}
Although Equation~\ref{eq:f1-macro} is an algebraic consequence of the
$F_1$ definition, its \emph{predictive} content is non-trivial: the
per-class weights $w^c$ vary across classes (empirically $0.37$--$0.66$
in our NER experiments), whereas the linear formula assumes a single
weight $\text{cov}$. The identity therefore explains \emph{why} the
linear formula's error is large ($\varepsilon \approx 0.019$) and
provides a constructive fix.


\subsection{Kronecker Task-Appropriateness}
\label{sec:kronecker}

A bilinear Kronecker layer computes $Y = A X B$ for $X \in \mathbb{R}^{n \times n}$.
Flattened, this is $Y_{\text{vec}} = (B^\top \otimes A) X_{\text{vec}}$,
a linear operator of size $n^2 \times n^2$ represented by only $2 n^2$
trainable parameters. The operator's rank is $\text{rank}(A) \cdot \text{rank}(B) = n^2$
(full rank), but its degrees of freedom remain $2 n^2$.

This structural constraint is a \emph{feature} when the task's own structure
aligns with the tensor decomposition. We formalise this intuition as:

\paragraph{Task-Appropriateness Hypothesis.}
Kronecker architectures are parameter-efficient when the task's functional
dependencies respect the tensor factorisation $X \mapsto A X B$. Conversely,
on tasks whose target is a generic (non-factorisable) linear map, the same
constraint becomes a bottleneck.

We test this hypothesis across three task types in Section~\ref{sec:experiments}.


\section{Experiments}
\label{sec:experiments}

\subsection{Linear Metric Validation}
\label{sec:exp-linear}

Equation~\ref{eq:linear} was validated on 12 datasets spanning 8 language
families, 4 writing systems (Latin, Han, Cyrillic, Arabic), and 4 task
types (binary dependency verification, multiclass, multi-label, NER).
Table~\ref{tab:linear} summarises the aggregate error.

\begin{table}[h]
\centering
\caption{Linear metric formula validation. Mean prediction error across
all configurations.}
\label{tab:linear}
\small
\begin{tabular}{lc}
\toprule
Dimension & Coverage \\
\midrule
Languages        & 12 (8 families) \\
Writing systems  & 4 (Latin, Han, Cyrillic, Arabic) \\
Task types       & 4 (binary, multiclass, multi-label, NER) \\
Metrics          & 2 (accuracy, bit-accuracy) \\
\midrule
Mean error       & $0.002$ \\
Max error        & $0.008$ \\
\bottomrule
\end{tabular}
\end{table}

Under symbolic noise (0--50\%), the formula's prediction error remains
invariant at $\varepsilon = 0.0066$, indicating that the linear identity
captures a structural property of the hybrid ensemble.


\subsection{F1 Per-Class Validation}
\label{sec:exp-f1}

The linear formula fails on $F_1$ because $F_1$ involves a harmonic mean
and is not linear in the per-item correctness indicator. We tested this
on 20 NER configurations (5 seeds $\times$ 4 hidden ratios).

\begin{table}[h]
\centering
\caption{$F_1$ formula comparison on 20 NER configurations. The linear
formula uses $\text{cov}$ as weight; the per-class formula uses
$w_c = D_s^c / (D_s^c + D_n^c)$.}
\label{tab:f1}
\small
\begin{tabular}{lcc}
\toprule
Formula & Mean $\varepsilon$ & Max $\varepsilon$ \\
\midrule
Linear (cov-weighted)     & $0.0188$ & $0.0515$ \\
Per-class ($w_c$-weighted) & $0.000000$ & $0.000000$ \\
\bottomrule
\end{tabular}
\end{table}

The per-class formula achieves machine-precision error on all 20
configurations, confirming that the extension is exact for macro $F_1$.


\subsection{Kronecker Task-Appropriateness Results}
\label{sec:exp-kronecker}

We tested the task-appropriateness hypothesis across three task types:
a random linear teacher (synthetic), a simple compositional task
(block powers), and an image-like compositional task
(2D grid with spatial correlation). Table~\ref{tab:kronecker} summarises
the winner in each case.

\begin{table}[h]
\centering
\caption{Winner by task type. Parameter counts in parentheses.
Statistical tests: paired $t$-test on matched seeds.}
\label{tab:kronecker}
\small
\begin{tabular}{lccc}
\toprule
Task type & Winner & Test Acc & Efficiency \\
\midrule
Random linear           & Dense MLP         & ---      & baseline \\
Simple compositional    & Flat Kronecker    & $0.8007$ & $8\times$ fewer \\
Image compositional     & Hier.\ Kronecker  & $0.8360$ & $34\times$ more \\
\bottomrule
\end{tabular}
\end{table}

\paragraph{Simple compositional.}
20 seeds. Flat Kronecker ($2{,}564$ params) reaches
$0.8007 \pm 0.0202$ versus Dense MLP ($20{,}868$ params) at
$0.7982 \pm 0.0190$; the difference is not significant
($p = 0.52$). Hierarchical Kronecker ($87{,}044$ params) reaches
$0.7872$ and is \emph{significantly worse} than flat
($p = 0.0008$, Cohen's $d = -0.75$).

\paragraph{Image compositional.}
50 seeds. Hierarchical Kronecker reaches $0.8360 \pm 0.0217$ versus
Dense MLP at $0.8190 \pm 0.0198$ (paired $t$-test:
$p < 0.0001$, Cohen's $d = -1.37$). Flat Kronecker
($2{,}564$ params) reaches $0.8267$ and is significantly better than
Dense MLP ($p = 0.0003$, $d = +0.51$), demonstrating $8\times$ parameter
efficiency with a small absolute accuracy advantage.

The contrast between the two compositional tasks---simple blocks versus
image-like grids with spatial correlation---demonstrates that
Kronecker's advantage depends on the \emph{degree} of structural alignment
between the task and the tensor factorisation.


\subsection{Real Image Validation (MNIST)}
\label{sec:mnist}

To test whether the task-appropriateness hypothesis extends beyond
synthetic grids, we evaluated on MNIST (28$\times$28 grayscale digits).
We trained on 2000 examples with 500 test examples for 30 epochs.

\begin{table}[h]
\centering
\caption{MNIST results (5 seeds). Flat Kronecker achieves competitive
accuracy with $12.6\times$ fewer parameters than the largest Dense MLP.}
\label{tab:mnist}
\small
\begin{tabular}{lcc}
\toprule
Model & Parameters & Test Accuracy \\
\midrule
Dense MLP ($h=32$)         & $26{,}506$  & $0.594 \pm 0.056$ \\
Dense MLP ($h=128$)        & $118{,}282$ & $0.808 \pm 0.017$ \\
Flat Kronecker ($K=1$)     & $9{,}418$   & $0.770 \pm 0.013$ \\
Flat Kronecker ($K=2$)     & $10{,}986$  & $0.717 \pm 0.041$ \\
Flat Kronecker ($K=3$)     & $12{,}554$  & $0.662 \pm 0.048$ \\
\bottomrule
\end{tabular}
\end{table}

\paragraph{Findings.}
(1)~\emph{Parameter efficiency extends to real images:} Flat
Kronecker ($9{,}418$ params) reaches $0.770$, versus Dense MLP
($118{,}282$ params) at $0.808$ --- a $12.6\times$ reduction with a
$3.8\%$ accuracy gap.
(2)~\emph{Depth collapse is universal:} on both synthetic and real
data, $K=1$ outperforms $K>1$, confirming that the collapse is
structural, not dataset-specific.
(3)~\emph{Structural constraint is regularising:} Dense MLP with
$h=32$ ($26{,}506$ params) underperforms Flat Kronecker with
$9{,}418$ params, suggesting that the Kronecker constraint acts as
an implicit regulariser.

\paragraph{Honest limitation.}
The absolute accuracy gap ($3.8\%$) remains. Kronecker is efficient,
not dominant.

\subsection{Negative Results}
\label{sec:negatives}

We document two negative results that bound the Kronecker claim.

\paragraph{Parameter-matched comparison.}
On the random linear task, when both architectures are given the same
parameter budget ($\sim 21{,}764$), Dense MLP reaches $0.900$ while
Flat Kronecker reaches $0.257$ at depth $K=168$. Deep flat chains
collapse: repeated bilinear + SiLU without normalisation destroys
information. Hierarchical Kronecker partially avoids collapse
($0.816$) but still loses to Dense MLP.

\paragraph{LayerNorm + residual.}
Adding LayerNorm and residual connections to flat Kronecker recovers
from $0.402$ to $0.584$, but does not close the gap to hierarchical
($0.817$) or dense ($0.916$). We conclude that the flat chain's
limitation is not purely a normalisation problem.

These results establish that the Kronecker advantage is
\emph{task-conditional}, not universal.


\section{Discussion}
\label{sec:discussion}

\subsection{Statistical Methodology}
\label{sec:stats}

All $p$-values reported in this paper are from paired $t$-tests on
matched seeds. Where multiple hypotheses are tested on the same task
(three pairwise comparisons per task), Bonferroni correction at
$\alpha = 0.05$ requires $p < 0.017$ for significance. Under this
criterion: flat vs.\ dense on image compositional
($p = 0.0003$), flat vs.\ hierarchical ($p < 0.0001$), and dense
vs.\ hierarchical ($p < 0.0001$) all remain significant. The
simple compositional task has one non-significant comparison
($p = 0.52$) and one significant (flat vs.\ hierarchical,
$p = 0.0008$).

\subsection{When to Use Which}
\label{sec:decision}

Our experiments yield a simple decision rule:

\begin{enumerate}
\item \textbf{Metric choice.} Use linear metrics (accuracy,
bit-accuracy) when possible; the gain formula is closed-form
(Equation~\ref{eq:linear}). For macro $F_1$, use the per-class
weighted identity.
\item \textbf{Architecture choice.} Match the architecture to the task
structure. Random linear tasks favour Dense MLP; block-structured
tasks favour Flat Kronecker; image-like tasks with spatial correlation
favour Hierarchical Kronecker.
\item \textbf{Hierarchy.} Do not deploy hierarchical Kronecker on
tasks without spatial structure: it costs $34\times$ more parameters
and underperforms.
\end{enumerate}

\subsection{Why Does the F1 Formula Require Per-Class Weights?}

For linear metrics, the weight collapses to coverage:
$\sum_c D_c = 2N$ identically, so $w = \text{cov}$.
For macro $F_1$, per-class weights $w_c = D_s^c / (D_s^c + D_n^c)$
differ across classes and from $\text{cov}$; the correct aggregation
is necessarily per-class.

\subsection{The Structural Alignment Principle}

The Kronecker results suggest a broader principle: \emph{a parameter
budget is most efficiently spent when its implicit constraints align
with the task's dependencies}. A bilinear layer constrains the operator
to the form $B^\top \otimes A$. When the task's target respects this
form, the constraint is free; when it does not, the constraint becomes
a bottleneck.


\section{Limitations and Future Work}
\label{sec:limitations}

\paragraph{Limitations.}
The Kronecker task-appropriateness study covers three synthetic task
types; generalisation to real-world image or language data is an
open question. The per-class $F_1$ identity is proven for macro $F_1$;
weighted and micro variants are not addressed. Our neural baselines
are small models trained from scratch; we do not compare against
pretrained language models. Statistical power is adequate
($n = 20$--$50$) but not exhaustive.

\paragraph{Future work.}
(1) Test the Kronecker hypothesis on natural images and text.
(2) Derive per-class identities for weighted $F_1$, precision, and recall.
(3) Extend the hybrid formula to pretrained large models.
(4) Investigate whether task-appropriateness can be predicted from
data statistics alone.

\section{Conclusion}
\label{sec:conclusion}

Hybrid systems do not win unconditionally. Their gain depends on three
independent dimensions: metric type (linear vs non-linear), architecture
type (Kronecker vs dense), and task structure (random vs compositional).
We provide a formula for each dimension, statistical evidence for the
architecture choice, and an honest account of when the geometric
architecture loses. Together, these results suggest that hybrid system
design can be made principled rather than heuristic.

\bibliographystyle{plain}
\begin{thebibliography}{9}

\bibitem{esa2026preprint1}
E.~Esa. Hybrid symbolic-neural paradigm: A universal formula for
predicting gain over neural baselines. \emph{arXiv preprint}, 2026.

\bibitem{garcez2023neurosymbolic}
A.~d'Avila Garcez and L.~C.~Lamb. Neurosymbolic AI: The 3rd wave.
\emph{Artificial Intelligence Review}, 2023.

\bibitem{manhaeve2018deepproblog}
R.~Manhaeve et al. DeepProbLog: Neural probabilistic logic programming.
\emph{NeurIPS}, 2018.

\bibitem{nivre2020ud}
J.~Nivre et al. Universal Dependencies v2. \emph{LREC}, 2020.

\bibitem{vaswani2017attention}
A.~Vaswani et al. Attention is all you need. \emph{NeurIPS}, 2017.

\end{thebibliography}

\end{document}
